\chapter*{Conventions Used in This Book}
\addcontentsline{toc}{chapter}{Conventions Used in This Book}

\section*{Language}

The text uses British English. Accordingly, terms such as \emph{colour} and
\emph{greyscale} are preferred in prose. Library, API and file-format names are
written exactly as their respective projects define them.

\section*{Code}

Python source code is shown in monospaced listings. Longer examples may be
loaded directly from the companion project so that the printed listing and the
working source remain synchronised.

\section*{Shell commands}

Commands intended for a terminal are shown without assuming a particular shell
unless the command itself is platform-specific. Chapter~2 will establish the
Anaconda, Visual Studio Code and Jupyter workflow used throughout the book.

\section*{Pixel coordinates}

Unless otherwise stated, image arrays use row-major NumPy indexing. An element
written as \texttt{image[y, x]} therefore identifies row \texttt{y} and column
\texttt{x}. Where Cartesian coordinates are used in mathematics, the coordinate
convention will be stated explicitly.

\section*{Historical weighting}

Asterisks such as \historicalstars{****} reproduce the historical relative
weighting of portfolio/sign-off functions where a source is available. They do
not indicate the importance of a topic in modern image processing.

\section*{References}

Numbered citations are managed with BibLaTeX and Biber. Historical course
material is cited for provenance. Technical descriptions should, wherever
possible, cite the relevant specification, original paper, authoritative text
or official software documentation.
